% =====================================================================
\chapter{Diseño de experimentos}
\label{tema:2}
% =====================================================================

Antes de que existiera la optimización bayesiana, la estadística industrial ya había resuelto una parte esencial del problema: decidir qué ensayos hacer \emph{antes} de ver ningún resultado para extraer el máximo de información con el mínimo de experimentos. Este tema sigue los capítulos~5, 6, 8 y~11 de Montgomery~\cite{montgomery}, texto de referencia de los cursos de diseño de experimentos, y las lecciones~6 y~11 del curso STAT~503 de Penn State~\cite{psu503}. Las superficies de respuesta proceden de Box y Wilson~\cite{boxwilson}; el hipercubo latino, de McKay, Beckman y Conover~\cite{mckay}, y los diseños para experimentos de ordenador, del capítulo~13 de Kochenderfer y Wheeler~\cite{kw}.

\section{Principios y variables codificadas}

Montgomery~\cite[§1.4]{montgomery} resume el diseño de experimentos en tres principios. La \emph{aleatorización} del orden de los ensayos protege contra factores no controlados que derivan en el tiempo (calentamiento de la máquina, humedad); la \emph{replicación} permite estimar el error puro y promediarlo; el \emph{bloqueo} agrupa los ensayos en condiciones homogéneas (un mismo lote de materia prima, un mismo turno) para eliminar esa fuente de variabilidad de la comparación.

\begin{definicion}[Variables codificadas]\label{def:codificadas}
Si la variable física $x_j$ se ensaya en el intervalo $[a_j,b_j]$, su versión codificada es
\[
u_j=\frac{x_j-c_j}{h_j},\qquad c_j=\frac{a_j+b_j}{2},\quad h_j=\frac{b_j-a_j}{2},
\]
de modo que el nivel bajo corresponde a $u_j=-1$, el alto a $u_j=+1$ y el centro a $u_j=0$.
\end{definicion}

Con variables codificadas todos los coeficientes de un modelo tienen las mismas unidades que la respuesta y son comparables entre sí; además, las columnas de los diseños factoriales resultan ortogonales~\cite[§6.9]{montgomery}.

\section{El diseño factorial \texorpdfstring{$2^k$}{2\^{}k}}

\begin{definicion}[Factorial $2^k$ y efectos]\label{def:factorial}
Un diseño factorial $2^k$ ensaya las $N=2^k$ combinaciones de $k$ factores a dos niveles $u_j\in\{-1,+1\}$, cada una con $n$ réplicas. Su \emph{matriz de diseño} $U\in\{-1,1\}^{N\times k}$ se escribe en \emph{orden estándar} (la primera columna alterna en cada fila, la segunda cada dos, la $j$-ésima cada $2^{j-1}$). La columna de la interacción entre los factores de un conjunto $S$ es el producto elemento a elemento $\prod_{j\in S}U_{\cdot j}$. El \emph{efecto} del término asociado a una columna $m\in\{-1,1\}^{N}$ es
\[
\hat E=\bar y_{+}-\bar y_{-}=\frac{2}{N}\sum_{i=1}^{N}m_i\,\bar y_i,
\]
donde $\bar y_i$ es la media de las réplicas del ensayo $i$ y $\bar y_{\pm}$ son las medias de los ensayos con $m_i=\pm1$. La cantidad $\sum_i m_i\,(n\bar y_i)$ se llama \emph{contraste}.
\end{definicion}

\begin{proposicion}[Ortogonalidad]\label{prop:ortogonal}
Sea $M\in\{-1,1\}^{N\times(N-1)}$ la matriz con las columnas de todos los efectos principales e interacciones de un $2^k$, y $X=[\mathbf 1\ M]$. Entonces $X^\top X=N\,I_N$.
\end{proposicion}

\begin{corolario}\label{cor:coeficientes}
El ajuste por mínimos cuadrados del modelo $y=\beta_0+\sum_{S}\beta_S\prod_{j\in S}u_j+\varepsilon$ a las medias $\bar y_i$ tiene solución $\hat\beta_S=\frac1N\sum_i m_{iS}\bar y_i=\hat E_S/2$: cada coeficiente es la mitad del efecto, y no depende de qué otros términos se incluyan en el modelo.
\end{corolario}

\begin{teorema}[Optimalidad del $2^k$]\label{teo:doptimo}
Sea un diseño cualquiera de $N$ ensayos en $[-1,1]^{k}$ para el modelo de primer orden $y=\beta_0+\sum_j\beta_ju_j+\varepsilon$, con errores independientes de varianza $\sigma^2$ y matriz $X$ de rango completo. Entonces $\Var(\hat\beta_j)\ge\sigma^2/N$ para todo $j$, y el factorial $2^k$ (o cualquier fracción ortogonal de él con $N$ ensayos) alcanza la cota.
\end{teorema}

Montgomery~\cite[§6.7]{montgomery} enuncia este resultado diciendo que los $2^k$ son diseños óptimos: con el mismo número de ensayos, ningún otro estima los efectos con menor varianza.

\section{Análisis de la varianza}

El análisis de la varianza (ANOVA) decide qué efectos son distinguibles del ruido. La teoría de contrastes es la del módulo~0.9, §5; aquí se necesita la versión para efectos de un factorial.

\begin{proposicion}[Descomposición de la variabilidad]\label{prop:anova}
En un $2^k$ con $n$ réplicas, sea $SS_T=\sum_{i,r}(y_{ir}-\bar y)^2$ la suma de cuadrados total. Entonces
\[
SS_T=\sum_{S\ne\varnothing}SS_S+SS_E,\qquad SS_S=\frac{nN}{4}\hat E_S^{2},\qquad SS_E=\sum_{i,r}(y_{ir}-\bar y_i)^2,
\]
con $N-1$ grados de libertad para los efectos (uno cada uno) y $N(n-1)$ para el error. Si los errores son normales de varianza $\sigma^2$ y $E_S=0$, el cociente $F_S=SS_S/MS_E$, con $MS_E=SS_E/(N(n-1))$, sigue una distribución $F(1,N(n-1))$.
\end{proposicion}

\section{Factoriales fraccionados}

Con $k=7$ factores, el factorial completo tiene 128 ensayos, y la mayoría de sus grados de libertad estiman interacciones de tres o más factores, casi siempre despreciables. Los factoriales fraccionados ensayan solo una fracción $2^{-p}$ del total~\cite[cap.~8]{montgomery}.

\begin{definicion}[Fracción $2^{k-p}$, alias y resolución]\label{def:fraccion}
Una fracción $2^{k-p}$ se construye con un factorial completo en $k-p$ factores básicos y $p$ \emph{generadores} que definen los demás como productos de los básicos (por ejemplo $D=ABC$). Cada generador da una palabra de la \emph{relación definidora} ($I=ABCD$); el conjunto de todas sus palabras y productos es el subgrupo definidor. Dos efectos son \emph{alias} si sus columnas coinciden en la fracción: el alias de un efecto se obtiene multiplicándolo por cada palabra y simplificando con $A^2=I$. La \emph{resolución} es la longitud de la palabra más corta de la relación definidora.
\end{definicion}

En resolución III los efectos principales tienen alias en interacciones dobles; en resolución IV, los principales están libres de dobles pero las dobles tienen alias entre sí; en resolución V, principales y dobles están libres entre sí~\cite[§8.6–8.7]{montgomery}. La elección habitual en una fase de cribado es resolución IV.

\addtocounter{definicion}{1}%

\section{Superficies de respuesta}

Los factoriales a dos niveles estiman efectos lineales e interacciones, pero no curvatura: no pueden localizar un óptimo interior. La metodología de superficies de respuesta de Box y Wilson~\cite{boxwilson} procede por etapas: un modelo de primer orden para avanzar por la dirección de máximo descenso hasta la región del óptimo y, allí, un modelo de segundo orden~\cite[§11.1–11.3]{montgomery}, \cite[lección~11]{psu503}.

\begin{definicion}[Modelo de segundo orden y diseño central compuesto]\label{def:ccd}
El modelo de segundo orden en $k$ variables codificadas es
\[
\hat y(u)=b_0+b^\top u+u^\top B u,\qquad B=B^\top,\quad B_{jj}=b_{jj},\ B_{jl}=\tfrac12 b_{jl}\ (j\ne l),
\]
con $1+k+k(k+1)/2$ coeficientes. El \emph{diseño central compuesto} (CCD) combina un $2^k$ (o una fracción de resolución V), $2k$ puntos axiales $(\pm\alpha,0,\dots,0),\dots$ y $n_c$ puntos centrales. Es \emph{rotable} (la varianza de la predicción solo depende de $\norm{u}$) si $\alpha=(2^k)^{1/4}$.
\end{definicion}

\begin{proposicion}[Punto estacionario y análisis canónico]\label{prop:estacionario}
Si $B$ es invertible, el único punto crítico de $\hat y$ es $u_s=-\tfrac12B^{-1}b$, con valor $\hat y(u_s)=b_0+\tfrac12 b^\top u_s$. Si $B=Q\Lambda Q^\top$ es su descomposición espectral y $w=Q^\top(u-u_s)$, entonces $\hat y(u)=\hat y(u_s)+\sum_j\lambda_jw_j^2$. El punto es un mínimo si todos los $\lambda_j>0$, un máximo si todos son negativos y un punto de silla en otro caso.
\end{proposicion}

\addtocounter{definicion}{1}%

\section{Diseños para experimentos de ordenador: el hipercubo latino}

Cuando la respuesta procede de una simulación determinista no tiene sentido replicar, y cuando se va a ajustar un modelo flexible (\cref{tema:3}) interesa cubrir el dominio uniformemente en lugar de concentrar los ensayos en los vértices. Montgomery~\cite[§11.5]{montgomery} y Kochenderfer y Wheeler~\cite[§13.3–13.5]{kw} tratan estos diseños de relleno del espacio.

\begin{definicion}[Hipercubo latino]\label{def:lhs}
Un \emph{hipercubo latino} (LHS) de $n$ puntos en $[0,1)^d$ es un conjunto $\{z_1,\dots,z_n\}$ tal que, para cada coordenada $j$, cada estrato $[k/n,(k+1)/n)$, $k=0,\dots,n-1$, contiene exactamente una de las $z_{ij}$. Se construye con $d$ permutaciones independientes $\pi_j$ de $\{0,\dots,n-1\}$ y uniformes $v_{ij}\in[0,1)$: $z_{ij}=(\pi_j(i)+v_{ij})/n$.
\end{definicion}

Toda proyección de un LHS sobre una coordenada está perfectamente repartida: si al final solo una variable resulta relevante, los $n$ ensayos dan $n$ niveles distintos de ella, mientras que una rejilla de $n$ puntos en $d$ dimensiones da solo $n^{1/d}$. Este argumento reaparece en el \cref{tema:6} para la búsqueda aleatoria.

\begin{teorema}[Reducción de varianza en una dimensión]\label{teo:estratificado}
Sea $g\colon[0,1)\to\R$ de cuadrado integrable, $\mu=\int_0^1g$, y sean $\hat\mu_{MC}=\frac1n\sum_ig(V_i)$ con $V_i$ uniformes independientes y $\hat\mu_{E}=\frac1n\sum_{k}g\bigl((k+V_k)/n\bigr)$ el estimador estratificado (un LHS en $d=1$). Ambos son insesgados y $\Var(\hat\mu_E)\le\Var(\hat\mu_{MC})$.
\end{teorema}

En dimensión $d$, McKay, Beckman y Conover~\cite[§2.1]{mckay} demostraron que el LHS reduce la varianza del estimador de la media respecto al muestreo aleatorio siempre que la respuesta sea monótona en cada variable. En la práctica, al estimar el coste medio de la inyectora en todo el dominio con 20 puntos, la desviación típica del estimador baja de $0{,}0247$ (Monte Carlo) a $0{,}0129$ (LHS): la varianza se reduce a la cuarta parte.

\begin{definicion}[Criterio maximín]\label{def:maximin}
Entre varios diseños de $n$ puntos se prefiere el que maximiza la distancia mínima entre pares, $\min_{i<l}\norm{z_i-z_l}$, o, equivalentemente en espíritu, el que minimiza una \emph{discrepancia} (diferencia entre la fracción de puntos en cada caja y su volumen).
\end{definicion}

Un LHS aleatorio puede ser malo (en el peor caso, todos los puntos en la diagonal); por eso se optimiza. En la práctica, con 10 puntos en tres dimensiones, la distancia mínima media pasa de $0{,}168$ con puntos uniformes a $0{,}228$ con un LHS, y la optimización \texttt{random-cd} de SciPy reduce la discrepancia centrada de $0{,}0154$ a $0{,}0105$.

\section{Caso de estudio: cribado de la inyectora con un \texorpdfstring{$2^3$}{2\^{}3}}

\addtocounter{definicion}{1}%

\begin{table}[ht]
\centering\small
\begin{tabularx}{\textwidth}{@{}>{\raggedright\arraybackslash}p{2.7cm}>{\raggedright\arraybackslash}X>{\raggedright\arraybackslash}p{2.1cm}>{\raggedright\arraybackslash}p{3.4cm}@{}}
\toprule
Herramienta & Para qué sirve & Licencia & Cuándo elegirla \\
\midrule
pyDOE3~\cite{pydoe3} & Factoriales completos y fraccionados, Plackett--Burman, CCD, Box--Behnken, LHS, diseños óptimos & BSD-3 & Generar diseños clásicos en Python \\
\texttt{scipy.\allowbreak stats.\allowbreak qmc}~\cite{scipyqmc} & LHS (con optimización \texttt{random-cd} y \texttt{lloyd}), Sobol, Halton, discrepancias & BSD & Diseños de relleno del espacio \\
statsmodels & ANOVA y regresión con tablas de contraste & BSD-3 & Análisis estadístico completo \\
R (FrF2, rsm) & Fraccionados con tablas de alias; superficies de respuesta & GPL & Estudios estadísticos detallados \\
Minitab, JMP & Diseño y análisis con interfaz gráfica & Comerciales & Entornos industriales con formación previa \\
\bottomrule
\end{tabularx}
\caption{Tecnologías para el diseño y el análisis de experimentos.}
\label{tab:tec2}
\end{table}

\addtocounter{definicion}{1}%

